\documentclass[11pt,a4paper]{article}
\usepackage{gtu-el-yazisi}
\begin{document}
\notbilgi{Makine Öğrenmesi}{Perceptron, backpropagation, aktivasyon ve PyTorch}{Ders-8.pdf}

\sayfa{1}
\section{ROC curve \okunamadi{}}
\elcizimi[0.4]{sekiller/Ders-8-s1-1.png}[1]
\[ \vec{w}\cdot\vec{x} + b = \hat{r}^t \]
\[ \mathcal{L}(\theta|X) = \sum_{t=1}^{N} (\hat{r}^t - r^t)^2 \]
türevini alıp 0'a \okunamadi{eşitleriz}

başka yöntem

\elcizimi[0.62]{sekiller/Ders-8-s1-2.png}[1]
sınıflandırmak için
\begin{gather*}
z<0 \Rightarrow C_0 \\
z\geq 0 \Rightarrow C_1 \\
P(C_0)
\end{gather*}

\subsection{sigmoid}
\[ \sigma(z) = \frac{1}{1+e^{-z}} \]
\elcizimi[0.4]{sekiller/Ders-8-s1-3.png}[1]
\begin{gather*}
z<0 \Rightarrow C_0 \\
z\geq 0 \Rightarrow C_1
\end{gather*}

Bu sefer iteratif olarak ilerliyoruz

gradient descent
\[ w_i \leftarrow w_i - \eta \frac{\partial J}{\partial w_i} \quad\rightarrow\quad \text{maliyet fonksiyonu} \]
\elcizimi[0.4]{sekiller/Ders-8-s1-4.png}[1]

\sayfa{2}
\section{Back Propagation}
\[ \hat{y}_i = \sigma(z_i) = \frac{1}{1+e^{-z_i}} \quad , \quad z_i = w^T \cdot x_i + b \]
\[ \text{entropy} \quad \sum_{i=1}^{k} p_i \log p_i \]
\[ J(w) = -\frac{1}{N}\sum_{i=1}^{N}\underbrace{y_i\log(\hat{y}_i) + (1-y_i)\log(1-\hat{y}_i)}_{\text{cross entropy}} \]
\[ \frac{\partial J(w)}{\partial w} = \underset{\text{\small\textcolor{turuncukalem}{\okunamadi{önce bu}}}}{\frac{\partial J(w)}{\partial \hat{y}_i}} \cdot \underset{\text{\small\textcolor{turuncukalem}{sonra bu}}}{\frac{\partial \hat{y}_i}{\partial z_i}} \cdot \underset{\text{\small\textcolor{turuncukalem}{sonra bu}}}{\frac{\partial z_i}{\partial w}} \]
\[ \frac{\partial J(w)}{\partial \hat{y}_i} = -\frac{y_i}{\hat{y}_i} + \frac{1-y_i}{1-\hat{y}_i} \qquad\qquad \frac{\partial \hat{y}_i}{\partial z_i} = \hat{y}_i(1-\hat{y}_i) \]
\[ \frac{\partial z_i}{\partial w} = x_i \]
\[ \frac{\partial J(w)}{\partial w} = \frac{1}{N}\sum_{i=1}^{N}(\hat{y}_i - y_i)\,x \]
\textcolor{turuncukalem}{sonra bunu ekleyeceğiz}\\
\textcolor{turuncukalem}{geri sararak yapıyoruz}
\[ \frac{\partial J(w)}{\partial b} = \frac{1}{N}\sum_{i=1}^{N}(\hat{y}_i - y_i) \]
\[ X = \{x^t, r^t\}_{t=1}^{M} \qquad P(C_0)=\hat{y}_i \qquad P(C_1)=(1-\hat{y}_i) \]
\elcizimi[0.7]{sekiller/Ders-8-s2-1.png}[2]

multi layer perceptron
\elcizimi[0.35]{sekiller/Ders-8-s2-2.png}[2]

activation function
\[ \phi(wx+b) \]

\sayfa{3}
\[ \Phi_2\Bigl(W \sum \Phi_1(wx+b)\Bigr) \]
Fourier'de
\[ \overset{\text{Basis function}}{\langle \Phi_1(x), \Phi_2(x)\rangle} = \int_{-\infty}^{\infty} \Phi_1(x)\,\Phi_2(x)\,dx = \overset{\text{ortogonal}}{0} \]
ama burda

\elcizimi[0.75]{sekiller/Ders-8-s3-1.png}[3]

Biz ortogonal olmalarını bekliyoruz

Böyle bir teorem var
\[ \sum_{i=1}^{\infty} w_i\,\Phi(x-x_i) \approx g(x) \]
yumuşak baz fonksiyonlarıyla mükemmel olarak tahmin edebiliriz

teorik olarak

universal approximator
\[ x \xrightarrow{f(x)} y \]
\elcizimi[0.45]{sekiller/Ders-8-s3-2.png}[3]

\[ \text{\textcircled{1}}\ \mathrm{softmax} = \frac{e^{z_i}}{\sum_{j=1}^{C} e^{z_j}} \quad \swarrow\ \text{olasılık dağılımına ihtiyacımız olduğunda} \]

\sayfa{4}
\okunamadi{soru}

\elcizimi[0.65]{sekiller/Ders-8-s4-1.png}[4]
böyle bir sıkıntı var \quad sonsuza giderken türev 0 \quad o yüzden back propagation iyi çalışmaz

bu yüzden genelde softmax en son katmanda kullanılır

\textcircled{2} ReLU

\elcizimi[0.35]{sekiller/Ders-8-s4-2.png}[4]
$ax+b$ şeklinde göstermiyoruz\\
non linear

negatif tarafta yine türevi 0, dead neuron sorununa yol açıyor

\textcircled{3} Leaky ReLU

\elcizimi[0.35]{sekiller/Ders-8-s4-3.png}[4]
nonlinearity daha az olduğu için daha fazla katman gerekiyor

\[ 180 \times [-1,1] \]
\textcircled{4} $\tanh = \dfrac{e^{z}-e^{-z}}{e^{z}+e^{-z}}$

\elcizimi[0.4]{sekiller/Ders-8-s4-4.png}[4]
bunda da türevler $+\infty,-\infty$'da 0'a yaklaşıyor

\sayfa{5}
\section{Derin öğrenmede katmanlar artar}
\textcircled{1} Tensorflow\\
\textcircled{2} Pytorch

TensorDataset

Data augmentation

\elcizimi[0.55]{sekiller/Ders-8-s5-1.png}[5]

Pseudo random

\okunamadi{x kütlesini} $\swarrow$\\
DataScaler() $\rightarrow$ Standart Scaler()
\[ \vec{x} = \begin{bmatrix} x_1 \\ x_2 \\ \vdots \\ x_N \end{bmatrix} \qquad -1 < x_i < 1 \]
\begin{gather*}
10^{6} \;\text{---}\; 10^{-4} \\
2 \;\text{---}\; 10^{-4}
\end{gather*}

\sayfa{6}
Scaler

\textcircled{1} min-max
\[ \frac{x-x_{\min}}{x_{\max}-x_{\min}} \]

\textcircled{2} Gauss $\mathcal{N}(0,1)$
\[ \frac{x-\mu}{\sigma} \]
\elcizimi[0.3]{sekiller/Ders-8-s6-1.png}[6]

sonra TensorDataset

random split

Dataloader

class ANNmodel(nn.Module) \quad $\leftarrow$ bunun inherit ediyoruz

init : super().\_\_init\_\_()

pytorch linear\\
torch.nn.Linear(in\_features, out\_features, bias=True, device=None, dtype=None)

\elcizimi[0.3]{sekiller/Ders-8-s6-2.png}[6]
\[ \text{parametre sayısı} = \underset{\text{in out}}{N\times M} + N \]

net = nn.Sequential

\hspace{2em}BatchNorm1d

\hspace{3em}Dropout

optim.Adam

nn.CrossEntropyLoss

\sayfa{7}
def forward(self,x):\\
\hspace*{2em}x = x.to(self.device)\\
\hspace*{2em}x = x.to(T.float32)

out\_ = F.one\_hot(T.argmax(out,dim=1), num\_classes=\okunamadi{k})

def learn:

self.optimizer.zero\_grad()\\
\hspace*{1em}y\_pred = self.forward(x)\\
\hspace*{2em}loss = self.loss(y\_pred, y)\\
\hspace*{3em}loss.backward

Normalization layer

\hspace{2em}BatchNorm1d

Dropout layers

overfit etmemek için bazı verileri düşürürüz

Cosine similarity

\elcizimi[0.3]{sekiller/Ders-8-s7-1.png}[7]
\[ \frac{x\cdot y}{\|x\|\,\|y\|} = \cos(\theta) \]

\sayfa{8}
Loss Functions

\hspace{2em}L1 loss

\hspace{2em}MSE loss

\hspace{2em}KLDivLoss

\end{document}
